% --- 题目 ---
\problem[P104-3 (16-3)]{设 $\displaystyle J_i = \iint_{D_i} \sqrt[3]{x-y} \, \mathrm{d}x \, \mathrm{d}y \ (i = 1,2,3)$, 其中
\[ D_1 = \{(x,y) \mid 0 \le x \le 1, 0 \le y \le 1\}, \]
\[ D_2 = \{(x,y) \mid 0 \le x \le 1, 0 \le y \le \sqrt{x}\}, \]
\[ D_3 = \{(x,y) \mid 0 \le x \le 1, x^2 \le y \le 1\}, \]
则 ( \quad )
\begin{tasks}(2)
\task $\displaystyle J_1 < J_2 < J_3$.
\task $\displaystyle J_3 < J_1 < J_2$.
\task $\displaystyle J_2 < J_3 < J_1$.
\task $\displaystyle J_2 < J_1 < J_3$.
\end{tasks}}
\ansat{301；【方法探究】 - 考点一：二重积分的概念与性质 - 3}

% --- 题目 ---
\problem[P104-2 (25-2)]{已知平面有界区域
\[ D = \left\{(x,y) \left| x^2 + y^2 \le 4x, x^2 + y^2 \le 4y \right.\right\}, \]
计算 $\displaystyle  \iint_D (x - y)^2 \, \mathrm{d}x \, \mathrm{d}y.$}
\ansat{301；【方法探究】 - 考点二：二重积分的计算 - 2}

% --- 题目 ---
\problem[P104-4 (24-1)]{已知平面区域
\[ D = \left\{(x,y) \left| \sqrt{1-y^2} \le x \le 1, -1 \le y \le 1 \right.\right\}, \]
计算 $\displaystyle \iint_D \frac{x}{\sqrt{x^2+y^2}} \mathrm{d}x \, \mathrm{d}y.$}
\ansat{301；【方法探究】 - 考点二：二重积分的计算 - 4}

% --- 题目 ---
\problem[P104-5 (24-2.3)]{设平面有界区域 $\displaystyle D$ 位于第一象限, 由曲线 $\displaystyle xy = \frac{1}{3}$, $\displaystyle xy = 3$ 与直线 $\displaystyle y = \frac{x}{3}, y = 3x$ 围成, 计算 $\displaystyle \iint_D (1+x-y) \mathrm{d}x \, \mathrm{d}y$.}
\ansat{302；【方法探究】 - 考点二：二重积分的计算 - 5}

% --- 题目 ---
\problem[P104-6 (23-2)]{设平面有界区域 $\displaystyle D$ 位于第一象限, 由曲线 $\displaystyle x^2 + y^2 - xy = 1, x^2 + y^2 - xy = 2$ 与直线 $\displaystyle y = \sqrt{3}x, y = 0$ 围成,
计算 $\displaystyle \iint_D \frac{1}{3x^2 + y^2} \mathrm{d}x \, \mathrm{d}y.$}
\ansat{302；【方法探究】 - 考点二：二重积分的计算 - 6}

% --- 题目 ---
\problem[P105-7 (23-3)]{已知平面区域 $$\displaystyle D = \left\{(x,y) \left| (x-1)^2 + y^2 \le 1 \right.\right\},$$
计算二重积分 $\displaystyle \iint_D \left| \sqrt{x^2+y^2} - 1 \right| \, \mathrm{d}x \, \mathrm{d}y.$}
\ansat{302；【方法探究】 - 考点二：二重积分的计算 - 7}

% --- 题目 ---
\problem[P105-8 (22-1,2,3)]{已知平面区域
\[ D = \left\{(x,y) \left| y-2 \le x \le \sqrt{4-y^2}, 0 \le y \le 2 \right.\right\}, \]
计算 $\displaystyle I = \iint_D \frac{(x-y)^2}{x^2+y^2} \, \mathrm{d}x \, \mathrm{d}y.$}
\ansat{302；【方法探究】 - 考点二：二重积分的计算 - 8}

% --- 题目 ---
\problem[P105-9 (21-2)]{设平面区域 $\displaystyle D$ 由曲线 $\displaystyle \left( x^2 + y^2 \right)^2 = x^2 - y^2$ ($\displaystyle x \ge 0, y \ge 0$) 与 $\displaystyle x$ 轴围成, 计算二重积分 $\displaystyle \iint_D xy \, \mathrm{d}x \, \mathrm{d}y.$}
\ansat{303；【方法探究】 - 考点二：二重积分的计算 - 9}

% --- 题目 ---
\problem[P105-10 (21-3)]{设有界区域 $\displaystyle D$ 是圆 $\displaystyle x^2 + y^2 = 1$ 和直线 $\displaystyle y = x$ 以及 $\displaystyle x$ 轴在第一象限围成的部分, 计算二重积分
\[ \iint_D \mathrm{e}^{(x+y)^2} (x^2 - y^2) \, \mathrm{d}x \, \mathrm{d}y. \]}
\ansat{303；【方法探究】 - 考点二：二重积分的计算 - 10}

% --- 题目 ---
\problem[P105-11 (20-2)]{设平面区域 $\displaystyle D$ 由 $\displaystyle x=1, x=2, y=x$ 与 $\displaystyle x$ 轴围成,
计算
\[ \iint_D \frac{\sqrt{x^2+y^2}}{x} \mathrm{d}x \, \mathrm{d}y. \]}
\ansat{303；【方法探究】 - 考点二：二重积分的计算 - 11}

% --- 题目 ---
\problem[P105-13 (19-2)]{已知平面区域
\[ D = \left\{(x,y) \left| |x| \le y, (x^2 + y^2)^3 \le y^4 \right.\right\}, \]
计算二重积分
\[ \iint_D \frac{x+y}{\sqrt{x^2+y^2}} \, \mathrm{d}x \, \mathrm{d}y. \]}
\ansat{303；【方法探究】 - 考点二：二重积分的计算 - 13}

% --- 题目 ---
\problem[P106-14 (18-2)]{设平面区域 $\displaystyle D$ 由曲线
\[ \begin{cases} x = t - \sin t, \\ y = 1 - \cos t \end{cases} (0 \le t \le 2\pi) \]
与 $\displaystyle x$ 轴围成, 计算二重积分
\[ \iint_D (x + 2y) \mathrm{d}x \, \mathrm{d}y. \]}
\ansat{303；【方法探究】 - 考点二：二重积分的计算 - 14}

% --- 题目 ---
\problem[P106-15 (18-3)]{设平面区域 $\displaystyle D$ 由曲线 $\displaystyle y = \sqrt{3(1-x^2)}$ 与直线 $\displaystyle y = \sqrt{3}x$ 及 $\displaystyle y$ 轴围成, 计算二重积分
\[ \iint_D x^2 \, \mathrm{d}x \, \mathrm{d}y. \]}
\ansat{304；【方法探究】 - 考点二：二重积分的计算 - 15}

% --- 题目 ---
\problem[P106-17 (17-3)]{计算 $\displaystyle \iint_D \frac{y^3}{(1+x^2+y^4)^2} \mathrm{d}x \, \mathrm{d}y$, 其中 $\displaystyle D$ 是第一象限中以曲线 $\displaystyle y=\sqrt{x}$ 与 $\displaystyle x$ 轴为边界的无界区域.}
\ansat{304；【方法探究】 - 考点二：二重积分的计算 - 17}

% --- 题目 ---
\problem[P106-18 (16-1)]{已知平面区域
\[ D = \left\{(r,\theta) \left| 2 \le r \le 2(1 + \cos\theta), -\frac{\pi}{2} \le \theta \le \frac{\pi}{2} \right.\right\}, \]
计算二重积分
\[ \iint_D x \, \mathrm{d}x \, \mathrm{d}y. \]}
\ansat{304；【方法探究】 - 考点二：二重积分的计算 - 15}

% --- 题目 ---
\problem[P107-6 (20-2)]{$\displaystyle \int_0^1 \mathrm{d}y \int_{\sqrt{y}}^1 \sqrt{x^3 + 1} \, \mathrm{d}x = \underline{\hspace{4em}}. $}
\ansat{305；【方法探究】 - 考点三：二次积分的积分次序与坐标系的转换 - 6}

% --- 题目 ---
\problem[P107-8 (16-2)]{已知函数 $\displaystyle f(x)$ 在 $\displaystyle \left[0, \frac{3\pi}{2}\right]$ 上连续, 在 $\displaystyle \left(0, \frac{3\pi}{2}\right)$ 内是函数 $\displaystyle \frac{\cos x}{2x - 3\pi}$ 的一个原函数, 且 $\displaystyle f(0) = 0$.
\begin{enumerate}
\item[(1)] 求 $\displaystyle f(x)$ 在区间 $\displaystyle \left[0, \frac{3\pi}{2}\right]$ 上的平均值;
\item[(2)] 证明 $\displaystyle f(x)$ 在区间 $\displaystyle \left(0, \frac{3\pi}{2}\right)$ 内存在唯一零点.
\end{enumerate}}
\ansat{305；【方法探究】 - 考点三：二次积分的积分次序与坐标系的转换 - 8}
